\chapter{Introduction}
\thispagestyle{empty}

Advances in modern science are closely linked to the ability to perform precise measurements. Many scientific breakthroughs have only become possible because advances in experimental techniques enabled researchers to implement and verify new concepts with increasing accuracy.
Therefore precise positioning systems are essential in many experimental and optical
applications, where even small positioning errors can significantly influence
measurement accuracy and system performance \prismcite{1}. 

One of the most accurate and widely used methods for displacement and position measurements is optical interferometry \prismcite{5}. It is based on the interference of coherent light waves, where changes in the optical path length produce measurable shifts in the interference pattern. Since the wavelength of light is on the order of hundreds of nanometers, interferometric techniques can resolve displacements that are only a small fraction of a wavelength, enabling sub-nanometer or even picometer resolution under suitable conditions \prismcite{41}.

One particularly demanding application is field-resolved infrared spectroscopy (FRS) based on electro-optic sampling (EOS), which enables the direct measurement of molecular electric-field responses in the time domain. By recording the coherent response of vibrationally excited molecules, FRS provides highly specific molecular fingerprints of biological samples such as blood plasma, opening new possibilities for biomedical diagnostics, health monitoring, and early disease detection, including cancer screening \prismcite{3,4}.

In the experimental setup used for these measurements, a motorized translation stage introduces a controlled optical delay between a gate pulse and the measured infrared signal. Accurate knowledge of the stage position is crucial, as the reconstructed electric-field waveform and all subsequent spectral information directly depend on the temporal delay. The currently implemented control software, however, only provides limited positioning accuracy and does not allow independent verification of the actual displacement. 
To improve the positioning accuracy of the system, an interferometer-based
calibration method was developed in this work.
